\subsection{分次}\label{sec:18d}
本节余下部分始终假设 $E$ 带有二次复体积形式 $\eta_E^2$。除以底空间 $S$ 上的形式 $dz^2$，得到相对形式 $\eta_{E/S}^2=\eta_E^2/dz^2$。$\eta_E^2$ 到 $\widetilde E$ 的拉回在分歧集 $M=E_*$ 之外非零。为了解该分歧集上的情形，写 $\eta_E^2=\eta_{E/S}^2\wedge dz^2$；由于该纤维非奇异，这在 $M$ 附近有意义。拉回后出现因子 $4y^2dy^2$，因此可在 $\widetilde E$ 上定义二次复体积形式及相应的相对形式，令
\begin{equation}\label{eq:18.15}
\eta_{\widetilde E}^2=\frac{\beta_{\widetilde E}^*\eta_E^2}{4y^2},\qquad
\eta_{\widetilde E/\widetilde S}^2=\eta_{\widetilde E}^2/dy^2.
\end{equation}
注意，这样安排后，非奇异纤维 $\widetilde E_y$ 上诱导的形式，与 $E_z$ 上的形式一致，其中 $z=\beta_{\widetilde S}(y)$。\znote{275}{因子 $4y^2$ 就是换元的 Jacobi 因子：由 $z=y^2+*$ 得 $dz=2y\,dy$，故 $dz^2=4y^2dy^2$；除以它之后 $\eta^2_{\widetilde E}$ 在底坐标 $z$ 下就是 $\beta_{\widetilde E}^*\eta^2_{E/S}\wedge dy^2$。若不作这一步，拉回的 $\beta_{\widetilde E}^*\eta^2_E$ 在分歧集 $M=\{y=0\}$ 上会带上二阶零点，这正是本节开头说「拉回在 $M$ 之外非零」的原因。又注意 $\eta_{\widetilde E}=\beta_{\widetilde E}^*\eta_E/2y$ 的符号依赖分支的选取，笔记 18.8 关于两个分支定向相反的讨论就建立在这一选择上。[来源：原书 §18d（印刷页 275）]}同样的论证可以应用于 $\Efamily$，将其看成沿 $\{w=\pi(x)\}$ 分歧的 $\mathcal W\times E$ 的二重覆盖。其结果
\SourcePage{283}
是，Lefschetz 纤维化 $\Psi:\Efamily\to\mathcal W$ 带有形式
\[
\eta_{\Efamily}^2=\frac{\beta_{\Efamily}^*\eta_E^2\wedge dw^2}{4y^2},\qquad
\eta_{\Efamily/\mathcal W}^2=\eta_{\Efamily}^2/dw^2,
\]
满足：在正则纤维 $\Efamily_w$ 上诱导的形式等于 $(\beta_{\Efamily}|\Efamily_w)^*\eta_E^2/4y^2$；特别地，$w=*$ 时就得到\eref{18.15}。回顾，在这种情形下，沿底空间 $\mathcal W$ 正则部分中任意环路的单值变换映射，自然成为分次辛自同构。

\StatementNumber{6}\begin{lemma}\label{st:18.6}
取\stref{引理}{18.1}中的单值变换，设 $\gamma^\#$ 为其自然分次。那么，对所有 $(y,x)\in\widetilde E^{\mathrm{shrink}}$，都有 $\gamma^\#(y,x)=-1$。
\end{lemma}
\begin{proof}
由 $\eta_{\Efamily/\mathcal W}^2$ 的定义及\eref{18.12}，在每个满足假设的点 $(y,x)$ 处，有
\begin{equation}\label{eq:18.16}
h_{\lambda|[0,t]}^*(\eta_{\Efamily/\mathcal W}^2|\Efamily_{\lambda(t)})
=\eta_{\widetilde E}^2\cdot\frac{\pi(x)-\lambda(0)}{\pi(x)-\lambda(t)}.
\end{equation}
当 $t$ 从 $0$ 变到 $1$ 时，\eref{18.16}中的分式绕原点顺时针转一周，因此\eref{15.8}中的所有相位函数也如此，所需结论随之成立。
\end{proof}

还须重新考察\sr{fig:18.1}{图~18.1}。以 $V_k,\Delta_k,\widetilde\Delta_k$ 表示与 $\gamma_k$ 相关的消失圈、Lefschetz 顶针及 (B) 型 Lagrangian 子流形；以 $V'_{k+1},\Delta'_{k+1},\widetilde\Delta'_{k+1}$ 表示从 $\gamma'_{k+1}$ 得到的相同数据。$V_k$ 与 $\gamma_{k+1}$ 的消失圈同痕，因而也与 $V'_{k+1}$ 同痕；更精确地说，用来构造 $\mu$ 的匹配圈的同痕 $I$，诱导了它们之间的同痕。由于 $\gamma_k$ 与 $\gamma'_{k+1}$ 在基点附近相同，这又诱导了 $\Delta_k\cap U$ 与 $\Delta'_{k+1}\cap U$ 之间的逐纤维同痕，其中 $U\subset E$ 为 $E_*$ 的一个适当小邻域。提升到二重覆盖后，在 $U$ 的原像 $\widetilde U$ 内，得到 $\widetilde\Delta_k\cap\widetilde U$ 与 $\widetilde\Delta'_{k+1}\cap\widetilde U$ 之间的同痕。选择与该局部同痕相容的分次，再将后一分次沿变形传递给 $\widetilde\Delta_{k+1}$。
\StatementNumber{7}\begin{lemma}\label{st:18.7}
为 $\widetilde\Delta_k$ 与 $\widetilde\Delta_{k+1}$ 如此选择分次，并为涉及的 Dehn 扭转赋予自然分次后，\stref{引理}{18.4}中的同痕成为分次 Lagrangian 子流形的同痕。
\end{lemma}
\begin{proof}
通常的变形技巧将问题归约到如下情形：匹配圈 $\Sigma_\mu$ 可以用朴素方式构造，且 $E$ 上的辛形式在 $\gamma_k(1)$ 上方的临界点附近，在局部全纯 Morse 坐标下为标准形式。加入第二个要求，是为了利用如下事实：在\stref{例题}{16.5}的局部模型中，沿底空间任意路径 $\beta$ 的平行移动，都会将消失圈 $\sqrt{\beta(0)}S^n$ 映到 $\sqrt{\beta(1)}S^n$。
\SourcePage{284}
仔细选取 $\gamma'_{k+1}$，使其在临界值 $\gamma_k(1)$ 的一个小邻域之外，与 $\gamma_k$ 或 $\mu$ 一致。此时，沿各平行移动映射追踪，便可看出 $V_k=V'_{k+1}$。因此，仍取上述 $\widetilde U$，有 $\widetilde\Delta_k\cap\widetilde U=\widetilde\Delta'_{k+1}\cap\widetilde U$。按定义，这两个子流形的分次相同。由于 Dehn 扭转集中在 $\widetilde\Sigma_{\mu,\pm}$ 附近，因而在 $\widetilde U$ 上平凡，故
\begin{equation}\label{eq:18.17}
\tau_{\widetilde\Sigma_{\mu,-}}\tau_{\widetilde\Sigma_{\mu,+}}(\widetilde\Delta_k)\cap\widetilde U
=\widetilde\Delta'_{k+1}\cap\widetilde U,
\end{equation}
而由 Dehn 扭转的典范分次诱导的分次仍然一致。现在，\stref{引理}{16.13}给出的\eref{18.17}中两个 Lagrangian 子流形之间的同痕，仍局限于 $\widetilde\Sigma_{\mu,\pm}$ 附近，因此分次相等的性质保持不变。
\end{proof}
